\section{Results and System Interpretation}
\subsection{STRICT Cross-Sample Parallelism}
On the prior eight-core GEMM cohort, STRICT Pipeline reduces construction from 149.56 to 20.95 seconds and achieves $1.638\pm0.005\times$ paired Process speedup over Stock; Group takes 37.88 seconds. Construction is about 45\% of Stock Process. Matched Group/Flat/Pipeline controls isolate C1 from other Stock repairs.

\begin{figure}[t]\centering\begin{tikzpicture}
\begin{groupplot}[group style={group size=2 by 1,horizontal sep=0.8cm},width=4.15cm,height=3.7cm,xtick={8,16},xmin=7,xmax=17,ymin=0,ymax=40,xlabel={Physical cores},ylabel={Build (s)},label style={font=\scriptsize},tick label style={font=\scriptsize},legend style={font=\scriptsize,at={(1.15,1.08)},anchor=south,legend columns=3},error bars/y dir=both,error bars/y explicit]
\nextgroupplot[title={},ylabel={Build (s)}]
\addplot+[mark=square*] coordinates {(8,17.413333333333334) +- (0,0.3250128202599597) (16,17.683333333333334) +- (0,0.3257811126098835)};
\addplot+[mark=triangle*] coordinates {(8,9.32) +- (0,0.15394804318340596) (16,5.506666666666667) +- (0,0.030550504633038766)};
\addplot+[mark=o] coordinates {(8,9.363333333333333) +- (0,0.10598742063723142) (16,5.3933333333333335) +- (0,0.08082903768654771)};
\legend{Group,Flat,Pipeline}
\nextgroupplot[title={},ylabel={}]
\addplot+[mark=square*] coordinates {(8,34.82333333333333) +- (0,0.3181718613160717) (16,33.99) +- (0,0.6602272336097627)};
\addplot+[mark=triangle*] coordinates {(8,19.553333333333335) +- (0,0.6217180497084931) (16,10.896666666666667) +- (0,0.040414518843273746)};
\addplot+[mark=o] coordinates {(8,18.993333333333332) +- (0,0.11503622617824955) (16,10.846666666666668) +- (0,0.04041451884327389)};
\end{groupplot}
\end{tikzpicture}

\caption{Construction means and sample SD, three processes per arm/point. Same-cohort Group/Flat/Pipeline controls on BNN (left) and DLA (right); capped routing tails do not establish full-route legality. Error bars are small at this scale.}\label{fig:causal}\end{figure}

In Fig.~\ref{fig:causal}, ten groups of four samples limit group-local concurrency. At eight cores, BNN Group/Flat/Pipeline take $17.41\pm0.33$, $9.32\pm0.15$ and $9.36\pm0.11$ s. Flat and Pipeline independently remove barriers and are nearly equal: the benefit comes from decoupling groups, not a particular Pipeline primitive. At sixteen cores Flat/Pipeline take 5.51/5.39 s, against Group's 17.68 s. DLA shows the same mechanism (34.82/19.55/18.99 s at eight cores). Separate 2D controls also benefit: eight-core BNN-2D takes 7.75/4.23/4.16 s. These are not equal-sized 2D/3D pairs, so layer count is not a causal explanation of time differences.

\subsection{Wave Scheduling: Concurrency Versus Work}
\begin{table}[t]\centering\caption{Matched BNN40 original-representation profiles; diagnostic operation counts, not native E2 timing.}\label{tab:waves}
\begin{tabular}{lrr}\toprule Metric & Shift 0 & Shift 20\\\midrule
\tableinput{tables/waves.tex}\bottomrule\end{tabular}\end{table}
Table~\ref{tab:waves} recomputes complete per-sample profiles. Shift 20 reduces summed waves by $9.112\times$ and maximum per-sample waves (the corresponding batch rounds) by $10.627\times$, while increasing eligible edge work only $0.3504\%$. Edge-work-weighted frontier grows $6.6484\times$. This exposes more work per synchronization opportunity; it is neither an occupancy measurement nor proof that launches explain all runtime. Labels remain FP32 and shift-0/20 complete CPU/CUDA outputs agree under ALT. Instrumented profile times are excluded from native statistics. Semantic fixture checks include eight random cases with native congestion differences, motivating separate semantic/QoR qualification.

\subsection{Matched Native Process and Construction}
\begin{table*}[t]\centering\caption{Complete native Process seconds and directly paired speedups, mean $\pm$ sample SD ($n=3$). S: Stock; P: Pipeline; C: ALT CUDA. GEMM/BNN are one new allocation; LSTM is a separate historical cohort.}\label{tab:matchedproc}
\begin{tabular}{lrrrrrr}\toprule Workload & Stock (s) & Pipeline (s) & ALT CUDA (s) & S/P & P/C & S/C\\\midrule
\tableinput{tables/matched_process.tex}\bottomrule\end{tabular}\end{table*}
\begin{table*}[t]\centering\caption{Native lookahead construction E2 seconds and directly paired speedups for the same runs as Table~\ref{tab:matchedproc}; mean $\pm$ sample SD, $n=3$.}\label{tab:matchede2}
\begin{tabular}{lrrrrrr}\toprule Workload & Stock (s) & Pipeline (s) & ALT CUDA (s) & S/P & P/C & S/C\\\midrule
\tableinput{tables/matched_e2.tex}\bottomrule\end{tabular}\end{table*}
CUDA/Pipeline E2 ratios are \GemmNewPipeBuild\ for GEMM and \BnnNewPipeBuild\ for BNN (Tables~\ref{tab:matchedproc}--\ref{tab:matchede2}). Their Process ratios, \GemmNewPipeProc\ and \BnnNewPipeProc, are smaller because Process retains loading, initialization, routing and checking. Ratios use matched rounds, never historical denominators. These implementation comparisons are not a complete two-factor ablation.

\subsection{Native LSTM: Five Matched Configurations}
\begin{table}[t]\centering\caption{Historical LSTM native seconds, mean $\pm$ sample SD ($n=3$); five arms on its own rotated eight-core/H100 cohort.}\label{tab:lstm}
\begin{tabular}{lrr}\toprule Method & E2 & Process\\\midrule
\tableinput{tables/lstm.tex}\bottomrule\end{tabular}\end{table}
CUDA obtains \LstmPipeBuild\ E2 and \LstmPipeProc\ Process acceleration against Pipeline; against Flat, \LstmFlatBuild\ and \LstmFlatProc; against matched ALT CPU, \LstmAltCpuBuild\ and \LstmAltCpuProc. ALT CPU is close to Pipeline: changed search semantics alone do not explain GPU acceleration. 

\subsection{Semantics, QoR and Size Dependence}
\begin{table}[t]\centering\caption{Native qualification: nodes/edges in millions; CPD/WL changes (\%) against Stock; $\Delta q$ counts final-map congestion-word changes against STRICT. All cases have two layers and 40 samples; LSTM is historical.}\label{tab:qor}
\begin{tabular}{lrrrrr}\toprule Case & $N/10^6$ & $M/10^6$ & CPD & WL & $\Delta q$\\\midrule
\tableinput{tables/qualification_compact.tex}\bottomrule\end{tabular}\end{table}
Complete ALT CPU/CUDA samples and maps agree on each input. Stock/Pipeline maps also agree, while ALT changes congestion representatives: 50,789 GEMM, 24,434 BNN and 23,111 LSTM words, with zero delay or missing-state differences. Native strict-delay ties and ALT whole-pair improvement remain distinct contracts. Each route independently passes FULL/STA with unchanged own route SHA and zero overuse; equivalent maps do not require identical route SHAs. Every new formal run passes the $+5\%$ CPD/WL gate against both paired controls; this is not a universal zero-loss guarantee.

The separate historical BNN Flat/CUDA cohort gives $3.274\pm0.222\times$ E2 and $1.091\pm0.011\times$ Process ratios; it supplies no denominator for the new Pipeline comparison. Its s38417 negative case remains $0.274\pm0.032\times$ E2 and $0.912\pm0.016\times$ Process. Device initialization, scans and launches can outweigh a small build. ALT CPU improves construction there, but Stock is fastest for Process; no automatic size selector or all-size benefit is established.

\subsection{Closest Alternative and Cache Applicability}\label{sec:applicability}
The locked executable exposes \texttt{compressed\_map}, whose embedded assertion and compute source require one layer; the evaluated native graphs have two. This is a static domain exclusion, not a timed failure or superiority claim over supported configurations. Its compressed reader and binary writer are unimplemented (CSV export is separate). Separable is exposed but not qualified or timed here. For ordinary map caching, 36 historical fixed-input reuse records have legal routing and independent FULL/STA checks. They establish valid cache hits, separate from fresh construction; no unobserved cross-design reuse is amortized.

\subsection{Resident Time Has a Different Scope}
A separate packed resident cohort excludes host preparation/full readback. Dual Gold 6530 uses 64 cores with a 55-CPU quota; 40 samples cap concurrency at 40. Among 8/32/64 workers, 64 gives the lowest means: $0.699624\pm0.026647$ s for BNN40 and $0.725510\pm0.021287$ s for LSTM40 (three processes, each averaging three repetitions after one warmup). H100 resident means are 0.151800/0.159498 s, including initialization, launches/control and completion synchronization. The BNN GPU record lacks an execution-time executable digest. These different representations/resources/scopes do not strengthen native E2/Process speedups.
